%% EC'26 camera-ready one-page abstract (HotCRP paper #2016).
%% Follows the camera-ready instructions emailed July 15, 2026 -- search
%% for "TODO" before uploading.
%%
%% Build: latexmk -pdf camera_ready.tex   (compiles clean; no -f needed)
%% Upload: https://ec2026.hotcrp.com/paper/2016/edit   (deadline July 29, 2026)
%% Template: ec26-proceedings-style-files.zip (pinned here: acmart.cls v2.12
%%   + acm-ec-26-proc.sty, both copied verbatim from the official zip)

\documentclass[format=acmsmall, review=false]{acmart}
\usepackage{acm-ec-26-proc}

%% Rights metadata: pasted verbatim from the ACM "Publication Release
%% Confirmation" email (rightsreview@hq.acm.org, July 27, 2026; a PDF copy
%% sits in this folder). License: CC BY 4.0. Conference line says *27th*.
\copyrightyear{2026}
\acmYear{2026}
\setcopyright{cc}
\setcctype{by}
\acmConference[EC '26]{The 27th ACM Conference on Economics and Computation}{July 06--10, 2026}{Rome, Italy}
\acmBooktitle{The 27th ACM Conference on Economics and Computation (EC '26), July 06--10, 2026, Rome, Italy}
\acmDOI{10.1145/3821539.3827858}
\acmISBN{979-8-4007-2813-6/2026/07}

\title[Chaining Tasks, Redefining Work]{Chaining Tasks, Redefining Work: A Theory of AI Automation}

%% Authors listed alphabetically; must match the HotCRP record (locked June 15, 2026).
%% ORCID iDs are omitted from the source for a uniform author block (not every author
%% has one registered); each author's ORCID still lives in their HotCRP profile.
\author{Mert Demirer}
\affiliation{%
  \institution{Massachusetts Institute of Technology}
  \city{Cambridge}
  \state{MA}
  \country{USA}}
\email{mdemirer@mit.edu}

\author{John J. Horton}
\affiliation{%
  \institution{Massachusetts Institute of Technology}
  \city{Cambridge}
  \state{MA}
  \country{USA}}
\affiliation{%
  \institution{NBER}
  \city{Cambridge}
  \state{MA}
  \country{USA}}
\email{jjhorton@mit.edu}

\author{Nicole Immorlica}
\affiliation{%
  \institution{Yale University}
  \city{New Haven}
  \state{CT}
  \country{USA}}
\affiliation{%
  \institution{Microsoft Research}
  \city{Cambridge}
  \state{MA}
  \country{USA}}
\email{nicimm@microsoft.com}

\author{Brendan Lucier}
\affiliation{%
  \institution{Microsoft Research}
  \city{Cambridge}
  \state{MA}
  \country{USA}}
\email{brlucier@microsoft.com}

\author{Peyman Shahidi}
\affiliation{%
  \institution{Massachusetts Institute of Technology}
  \city{Cambridge}
  \state{MA}
  \country{USA}}
\email{peymansh@mit.edu}

\begin{document}

%% Abstract: verbatim from the EC'26 submission (main.tex).
\begin{abstract}
We introduce a task-based model of production that responds to automation technologies such as AI.
%
Production is a sequence of \emph{steps} that firms aggregate into \emph{tasks} and then into \emph{jobs}, trading off gains from specialization against coordination costs.
%
Steps may be executed \emph{manually} by humans, \emph{augmented} with AI assistance under human oversight, or fully \emph{automated} within contiguous chains of AI-executed steps.
%
We characterize how firms optimally allocate steps between humans and AI in the short run, when job boundaries are fixed, and the long run, when AI deployment strategy and job design are jointly chosen.
%
We show analytically that job-level exposure to AI depends on the extent to which AI-exposed steps are adjacent or dispersed within the workflow.
%
Our framework also yields predictions about worker specialization and non-linearity of productivity gains from marginal improvements in AI quality.
%
Empirical evidence supports the model’s key predictions that: (1) AI-executed steps tend to co-occur in chains, (2) dispersion of AI-exposed steps lowers conversion from AI exposure to AI execution at the job level, and (3) adjacency to other AI-executed steps increases the likelihood of AI execution.
\end{abstract}

%% CCS codes generated per https://dl.acm.org/ccs -- at least one is required.
\begin{CCSXML}
<ccs2012>
   <concept>
       <concept_id>10010405.10010455.10010460</concept_id>
       <concept_desc>Applied computing~Economics</concept_desc>
       <concept_significance>500</concept_significance>
       </concept>
 </ccs2012>
\end{CCSXML}

\ccsdesc[500]{Applied computing~Economics}

%% Free-form keywords -- TODO: adjust wording with coauthors.
\keywords{AI Automation, Task Chaining, Division of Labor, Job Design}

\maketitle

\vspace{1cm}

\noindent
The full version of this paper is available at \url{https://arxiv.org/abs/2606.15960}.

%% No acknowledgements by design: not the norm in abstract-only entries.

\end{document}
